\documentclass[10pt]{article}
\usepackage[T1]{fontenc}
\usepackage[utf8]{inputenc}
\usepackage[french]{babel}
\usepackage{amsmath,amssymb,mathrsfs,mathpazo}
\usepackage{geometry}
\geometry{a4paper,margin=2cm}
\usepackage{graphicx,array,tabularx,tabularray,colortbl,multirow,multicol,booktabs,enumitem,pifont,ulem,fancybox,diagbox,cancel}
\usepackage{pstricks,pstricks-add,pst-plot,pst-tree,pst-node,pst-text,pst-func,pst-math,pst-bspline,pst-3dplot,pst-eucl,pst-solides3d}
\let\pscalculate\undefined
\usepackage{pst-bezier}
\usepackage{tikz,pgfplots,tkz-tab}
\usetikzlibrary{arrows,arrows.meta,patterns,calc,positioning,trees,decorations.markings,shapes.geometric,shapes.misc,shapes.arrows,intersections,decorations.pathmorphing,decorations.pathreplacing,backgrounds,shadows,babel}
\pgfplotsset{compat=1.18}
\usepgfplotslibrary{fillbetween}
\usepackage[np]{numprint}
\usepackage{listings,makecell,pas-tableur,twoopt}
\let\barre\undefined
\usepackage[pstricks]{tablvar}
\definecolor{darkgreen}{RGB}{0,100,0}
\providecommand{\red}{\color{red}}
\providecommand{\blue}{\color{blue}}
\providecommand{\green}{\color{green}}
\providecommand{\black}{\color{black}}
\providecommand{\white}{\color{white}}
\providecommand{\gray}{\color{gray}}
\providecommand{\cyan}{\color{cyan}}
\providecommand{\magenta}{\color{magenta}}
\providecommand{\yellow}{\color{yellow}}
\providecommand{\violet}{\color{violet}}
\providecommand{\orange}{\color{orange}}
\providecommand{\euro}{\texteuro}
\providecommand{\EUR}{\texteuro}
\providecommand{\ieme}{\textsuperscript{e}}
\providecommand{\degres}{\ensuremath{^\circ}}
\providecommand{\degree}{\ensuremath{^\circ}}
\providecommand{\vv}[1]{\overrightarrow{#1}}
\providecommand{\Coord}[1]{#1}
\providecommand{\decosix}{}
\providecommand{\starredbullet}{\ensuremath{\bullet}}
\pagestyle{empty}
\setlength{\parindent}{0pt}
\usepackage{adjustbox}\DeclareUnicodeCharacter{2236}{:}\DeclareUnicodeCharacter{2190}{\ensuremath{\leftarrow}}
\providecommand{\R}{\mathbb{R}}
\providecommand{\N}{\mathbb{N}}
\providecommand{\Z}{\mathbb{Z}}
\providecommand{\Q}{\mathbb{Q}}
\providecommand{\D}{\mathbb{D}}
\providecommand{\C}{\mathbb{C}}
\providecommand{\e}{{\mathrm{e}}}
\providecommand{\vect}[1]{\overrightarrow{#1}}
\providecommand{\vectt}[1]{\overrightarrow{\mathrm{#1}}}
\providecommand{\vvt}[1]{\overrightarrow{\mathrm{#1}}}
\providecommand{\np}[2][]{{#2~#1}}
\providecommand{\vv}[1]{{\overrightarrow{#1}}}
\renewcommand{\d}{{\mathrm{d}}}
\providecommand{\degres}{{^\circ}}
\providecommand{\degre}{{^\circ}}
\providecommand{\blue}{\color{blue}}
\providecommand{\red}{\color{red}}
\providecommand{\enskip}{\quad}
\providecommand{\llbracket}{\lbrack\!\lbrack}
\providecommand{\rrbracket}{\rbrack\!\rbrack}
\providecommand{\euro}{{\text{€}}}
\providecommand{\no}{{\text{n°}}}
\providecommand{\interval}[2]{{[#1;#2]}}
\providecommand{\intervFO}[2]{{[#1;#2[}}
\providecommand{\dotfill}{\ldots}

\providecommand{\ds}{\displaystyle}
\providecommand{\textding}[1]{#1}
\providecommand{\Oij}{\(\mathrm{O};\vec{\imath},\vec{\jmath}\)}
\providecommand{\Oijk}{\(\mathrm{O};\vec{\imath},\vec{\jmath},\vec{k}\)}
\providecommand{\Ouv}{\(\mathrm{O};\vec{u},\vec{v}\)}
\providecommand{\intervFF}[2]{\ensuremath{[#1;#2]}}
\renewcommand{\euro}{\ifmmode\text{\texteuro}\else\texteuro\fi}
\DeclareUnicodeCharacter{00B0}{\ifmmode^{\circ}\else\textdegree\fi}
\DeclareUnicodeCharacter{20AC}{\ifmmode\text{\texteuro}\else\texteuro\fi}
% Énoncés : Éducation nationale. Transcription et corrigés : APMEP, Denis Vergès et contributeurs. Découpage et classement : CoopMaths, Sébastien Lozano, Éric Elter et contributeurs. Adaptation de la présentation : Maths Explorer.
% CC BY-SA : https://coopmaths.fr/www/about/
% https://coopmaths.fr/alea/static/eam/2025/tex/eam_technologique_2025_06_sujet0v1_2.tex
\begin{document}
\noindent\textit{APMEP / CoopMaths. CC BY-SA. Présentation adaptée par Maths Explorer.}\par\medskip
Une biologiste désire étudier l'évolution de la population de singes sur une île. \\
En 2025, elle estime qu'il y a 1\,000 singes sur l'île.

\subsubsection*{A. Premier modèle}
Chaque année, la population de singes baisse de 10\%.

\begin{enumerate}
\item Montrer qu'en 2026, il y aura 900 singes sur l'île.

\item Pour tout entier naturel $n$, on note $u_n$ le nombre de singes sur l'île pour l'année $2025 + n$.

On a donc $u_0 = 1\,000$.

\begin{enumerate}
\item Indiquer ce que représente $u_2$ et calculer sa valeur.
\item Déterminer la nature de la suite $(u_n)$ et préciser sa raison.
\item Donner les variations de cette suite.
\end{enumerate}

\item Selon ce modèle, la population de singes est-elle menacée d'extinction ? Justifier.
\end{enumerate}

\bigskip

\begin{minipage}[t]{0.75\textwidth}
\subsubsection*{B. Second modèle}
On admet que l'évolution du nombre de singes est modélisée par la suite $(v_n)$ ainsi définie :

\[
\begin{cases}
v_{n+1} = 0,9v_n + 150~; & n \in \N \\
v_0 = 1\,000~,
\end{cases}
\]

où $v_n$ désigne le nombre de singes sur l'île pour l'année $2\,025 + n$.

\begin{enumerate}
\item Avec ce modèle, quelle sera la population de singes en 2\,026 ? \\
Détailler le calcul.

\item La feuille de calcul ci-contre donne les valeurs arrondies à l'unité des premiers termes de la suite $(v_n)$. \\
Quelle formule, destinée à être étirée vers le bas, faut-il saisir dans la cellule B3 pour obtenir les termes de la suite $(v_n)$ ?

\item Indiquer en quelle année la population de singes dépassera pour la première fois 1\,400 individus.
\end{enumerate}
\end{minipage}
\hfill
\begin{minipage}[t]{0.2\textwidth}
\vspace{0pt}
\centering
\renewcommand{\arraystretch}{1.2}
\begin{tabular}{|>{\columncolor{lightgray}}c|c|c|}
\hline
\rowcolor{lightgray}
\textbf{} & \textbf{A} & \textbf{B} \\ \hline
\textbf{1} & $n$ & $V_n$ \\ \hline
\textbf{2} & 0 & 1\,000 \\ \hline
\textbf{3} & 1 & 1\,050 \\ \hline
\textbf{4} & 2 & 1\,095 \\ \hline
\textbf{5} & 3 & 1\,136 \\ \hline
\textbf{6} & 4 & 1\,172 \\ \hline
\textbf{7} & 5 & 1\,205 \\ \hline
\textbf{8} & 6 & 1\,234 \\ \hline
\textbf{9} & 7 & 1\,261 \\ \hline
\textbf{10} & 8 & 1\,285 \\ \hline
\textbf{11} & 9 & 1\,306 \\ \hline
\textbf{12} & 10 & 1\,326 \\ \hline
\textbf{13} & 11 & 1\,343 \\ \hline
\textbf{14} & 12 & 1\,359 \\ \hline
\textbf{15} & 13 & 1\,373 \\ \hline
\textbf{16} & 14 & 1\,386 \\ \hline
\textbf{17} & 15 & 1\,397 \\ \hline
\textbf{18} & 16 & 1\,407 \\ \hline
\textbf{19} & 17 & 1\,417 \\ \hline
\textbf{20} & 18 & 1\,425 \\ \hline
\textbf{21} & 19 & 1\,432 \\ \hline
\end{tabular}
\end{minipage}

\end{document}
